\begin{proof}

\begin{enumerate}[label={$\langle1\rangle$\arabic*.},leftmargin=*]
\item The case $D>0$ is \ref{push:src:E.2.6}.  If $D<0$ and $E_0$ had a
horizontal prime component, the final part of \ref{push:src:E.2.8} would
give $b_1(F)=1$ for a smooth general fibre of $g_0$ and
would make every reduced fibre of $g_0$ irreducible.
\ref{push:src:E.2.9} would then forbid precisely that horizontal component.
Thus for either sign $E_0$ is vertical.

\item The connected-divisor-limit argument in the middle of the
proof of \ref{push:src:E.2.6} only uses verticality, not the sign of $D$.
Its generic choice of pencil section has central zero
support equal to the disjoint union of the vertical fixed
part and $e$ distinct regular fibres.  The total divisor
is reduced and irreducible with connected central support
after normalization and Stein factorization.  Consequently
$E_0=0$ and $e=1$.  Surjectivity of evaluation, flatness and
the torus assertion follow from the remainder of that same
proof.  The hypothesis of \ref{push:src:E.2.7} is now satisfied for either
sign, and gives the interior period limit.

\item Suppose two such one-sided families have the same central
manifold $X_0$.  Its algebraic reduction is intrinsic, so
both pencil lines are the same line
$L=g_0^*\mathcal O(1)$ up to isomorphism.  With indices
$i=1,2$, the two root identities give
\[
 K_{X_0}^{A_i}=L^{D_i},\qquad
 L^{D_1A_2-D_2A_1}=\mathcal O.
\]
The line $L$ is not torsion, because a positive power has
a nonzero section with zeros.  Hence $D_1/A_1=D_2/A_2$.
\SeriesRef{push:src:E.2.3} recovers the ordered coprime multiplicity pair from
this rational number.  This proves the common-limit
assertion without identifying the limit with a member of
the given period construction.

\end{enumerate}
\end{proof}
